\documentclass[11pt]{article}
\usepackage[a4paper,margin=22mm]{geometry}
\usepackage{fontspec}
\setmainfont{DejaVu Serif}
\setsansfont{DejaVu Sans}
\usepackage{microtype}
\usepackage{xcolor}
\usepackage{hyperref}
\usepackage{enumitem}
\usepackage{parskip}
\definecolor{Accent}{HTML}{971D32}
\definecolor{Muted}{HTML}{666666}
\hypersetup{colorlinks=true,urlcolor=Accent,linkcolor=Accent}
\setlist{leftmargin=1.6em,itemsep=0.3em,topsep=0.35em}
\setcounter{tocdepth}{2}
\renewcommand{\contentsname}{Contents}
\newcommand{\sectionrule}{\par\vspace{-0.5em}\noindent\textcolor{Accent}{\rule{\linewidth}{0.6pt}}\par}
\begin{document}

\begin{center}
{\sffamily\bfseries\color{Accent} TOEFL WORD OF THE DAY\par}
\vspace{8mm}
{\fontsize{42}{48}\selectfont\bfseries seep\par}
\vspace{2mm}
{\Large\sffamily\color{Accent} /siːp/\par}
\vspace{2mm}
{\small\sffamily Local audio phrase: seep\par}
{\small\sffamily Audio file: audios/seep.mp3\par}
\vspace{2mm}
{\large\sffamily\color{Muted} intransitive verb\par}
\vspace{5mm}
{\sffamily flow slowly through small openings \textbullet\ spread gradually\par}
\vspace{6mm}
{\large Seep describes liquid, gas, information, or influence moving slowly and gradually into, out of, or through something.\par}
\vspace{6mm}
\begin{minipage}{0.84\textwidth}
\itshape “Contaminated water began to seep through cracks in the storage tank.”
\end{minipage}\par
\vspace{10mm}
\textbf{Date:} September 13th 2026\quad\textbullet\quad \textbf{Level:} B1--mid-B2 TOEFL
\vfill
Created and compiled by \textbf{Amir Kooshky}\par
\vspace{2mm}
\href{https://t.me/KooshkyTOEFL}{Telegram Channel}\quad\textbullet\quad
\href{https://www.instagram.com/kooshkytoefl}{Instagram}\quad\textbullet\quad
\href{https://t.me/KooshkyTOEFL\_pv}{Personal Telegram}
\end{center}
\newpage
\tableofcontents
\newpage

\section{Meanings and usage}
\sectionrule
\subsection{Meaning 1: Flow slowly through small openings}
\textbf{Part of speech:} intransitive verb

\textbf{IPA:} /siːp/

\textbf{Pronunciation phrase:} seep

\textbf{Local audio file:} audios/seep.mp3

\textbf{Register:} neutral; common in environmental, scientific, and everyday descriptions

\textbf{Definition:} to flow or pass slowly in small amounts through tiny openings, a surface, or a material

\textbf{In simpler words:} leak or flow very slowly

\textbf{Typical pattern:} seep into, through, out of, from, or under + place or material

\subsubsection{Natural examples}
\begin{enumerate}
  \item Contaminated water began to seep through cracks in the storage tank.
  \item After several days of rain, moisture seeped into the basement walls.
  \item Oil from the damaged pipeline seeped into the surrounding soil.
  \item Groundwater can seep through layers of porous rock.
  \item A small amount of gas was seeping out of the damaged valve.
  \item Rainwater seeped under the door and formed a shallow pool on the floor.
  \item Chemicals from the landfill may seep into nearby wells.
  \item Sap slowly seeped from the cut in the tree trunk.
\end{enumerate}

\subsubsection{Nuance and implication}
\textbf{Central nuance:} Seep emphasizes slow movement in small quantities, often through pores, cracks, or another barrier.

\textbf{Usual implication:} The movement may be difficult to notice at first and may occur without deliberate control.

\textbf{Compared with pour:} Pour describes a much faster or more continuous flow. Seep describes slow movement in small amounts, often through tiny openings.

\subsubsection{Grammar notes}
\textbf{How to use it:} Seep does not normally take a direct object. Follow it with an adverb or a preposition that shows the direction of movement.

\textbf{What can follow it:} Common choices are into, through, out of, from, under, and across.

\textbf{Common preposition:} Use seep into for the destination, seep through for the material or opening crossed, and seep out of or from for the source.

\subsubsection{Useful patterns}
\textbf{Pattern:} seep into + place or substance

\textbf{Use:} Use this to name where the liquid or gas enters.

\textbf{Example:} Pesticides can seep into groundwater after heavy rain.

\textbf{Pattern:} seep through + material or opening

\textbf{Use:} Use this to name the barrier that the substance slowly crosses.

\textbf{Example:} Water seeped through a narrow crack in the concrete.

\textbf{Pattern:} seep out of or from + source

\textbf{Use:} Use this to identify where the substance slowly escapes.

\textbf{Example:} Gas was seeping out of a damaged pipe.

\textbf{Pattern:} seep under + barrier

\textbf{Use:} Use this when liquid passes slowly beneath something.

\textbf{Example:} Floodwater seeped under the temporary wall.

\subsubsection{Common wrong usage}
\paragraph{Correction 1}
\textbf{Wrong:} The crack seeped water into the room.

\textbf{Correct:} Water seeped into the room through the crack.

\textbf{Why:} Seep is normally intransitive: the liquid is the subject, not the object.

\paragraph{Correction 2}
\textbf{Wrong:} Water seeped the concrete wall.

\textbf{Correct:} Water seeped through the concrete wall.

\textbf{Why:} Use through to name the material that the liquid crosses.

\paragraph{Correction 3}
\textbf{Wrong:} The river seeped rapidly over its banks.

\textbf{Correct:} The river flowed rapidly over its banks.

\textbf{Why:} Seep implies slow movement in small amounts, so it does not fit a rapid large flow.

\subsection{Meaning 2: Spread or become known gradually}
\textbf{Part of speech:} intransitive verb

\textbf{IPA:} /siːp/

\textbf{Pronunciation phrase:} seep

\textbf{Local audio file:} audios/seep.mp3

\textbf{Register:} neutral; figurative and common in journalism and academic discussion

\textbf{Definition:} to spread, enter, escape, or become known slowly and gradually

\textbf{In simpler words:} spread or become understood little by little

\textbf{Typical pattern:} information or influence + seep into, out, through, or in

\subsubsection{Natural examples}
\begin{enumerate}
  \item News of the discovery gradually seeped out despite efforts to keep it confidential.
  \item The values of the movement seeped into mainstream political debate.
  \item A sense of uncertainty began to seep into the team's discussions.
  \item Ideas from the scientific revolution slowly seeped into popular culture.
  \item Details of the private negotiations eventually seeped out to the press.
  \item The stress of the conflict seeped into nearly every aspect of daily life.
  \item Foreign musical influences seeped into the region's traditional style.
  \item As the lecture continued, the significance of the evidence slowly seeped in.
\end{enumerate}

\subsubsection{Nuance and implication}
\textbf{Central nuance:} This figurative use presents information, emotion, or influence as if it were liquid moving slowly through a barrier.

\textbf{Usual implication:} The change is gradual, indirect, and sometimes difficult to prevent or identify.

\textbf{Compared with spread:} Spread is the general word and can describe any speed or method. Seep emphasizes a slow, gradual, and often unnoticed process.

\subsubsection{Grammar notes}
\textbf{How to use it:} The information, feeling, or influence is normally the subject. Add an adverb or preposition to show its direction.

\textbf{What can follow it:} Use into for an area being affected, out for information becoming public, or in for an idea gradually being understood.

\textbf{Common preposition:} Seep out often describes secret information becoming known; seep into describes an influence gradually affecting something.

\subsubsection{Useful patterns}
\textbf{Pattern:} information or news + seep out

\textbf{Use:} Use this when information becomes public slowly or indirectly.

\textbf{Example:} Reports of the disagreement began to seep out.

\textbf{Pattern:} idea, feeling, or influence + seep into + area

\textbf{Use:} Use this when something gradually begins to affect another area.

\textbf{Example:} Commercial influences seeped into the artistic movement.

\textbf{Pattern:} meaning or realization + seep in

\textbf{Use:} Use this when someone begins to understand something gradually.

\textbf{Example:} Only later did the importance of the warning seep in.

\subsubsection{Common wrong usage}
\paragraph{Correction 1}
\textbf{Wrong:} The journalist seeped the confidential news.

\textbf{Correct:} The confidential news seeped out to journalists.

\textbf{Why:} Seep is normally intransitive. The information that spreads is the subject.

\paragraph{Correction 2}
\textbf{Wrong:} The realization seeped me in.

\textbf{Correct:} The realization slowly seeped in.

\textbf{Why:} Do not place a person directly after seep in this pattern.

\section{Word family}
\sectionrule
\subsection{seepage}
\textbf{Part of speech:} uncountable noun

\textbf{IPA:} /ˈsiːpɪdʒ/

\textbf{Definition:} the slow movement of liquid or gas through small openings or a material

\textbf{Relationship to the headword:} This noun names the physical process or result of liquid or gas seeping.

\textbf{Grammar pattern:} seepage into, through, or from + place

\textbf{Usage note:} It is normally used without a or an and is mainly used for the physical meaning.

\subsubsection{Examples}
\begin{enumerate}
  \item Engineers detected water seepage beneath the dam.
  \item A protective layer prevents chemical seepage into the soil.
  \item The inspection revealed signs of seepage around the window.
\end{enumerate}

\section{Common collocations}
\sectionrule
\subsection{seep into the ground}
\textbf{Applies to:} Flow slowly through small openings

\textbf{Meaning:} pass slowly from the surface into the soil

\textbf{Example:} Some of the rainwater seeps into the ground and replenishes underground supplies.

\subsection{seep through cracks}
\textbf{Applies to:} Flow slowly through small openings

\textbf{Meaning:} pass slowly through narrow openings

\textbf{Example:} Water seeped through cracks in the foundation.

\subsection{seep into groundwater}
\textbf{Applies to:} Flow slowly through small openings

\textbf{Meaning:} move gradually through soil and contaminate underground water

\textbf{Example:} Agricultural chemicals may seep into groundwater.

\subsection{information seeps out}
\textbf{Applies to:} Spread or become known gradually

\textbf{Meaning:} information gradually becomes public

\textbf{Example:} Information about the negotiations slowly seeped out.

\subsection{seep into public awareness}
\textbf{Applies to:} Spread or become known gradually

\textbf{Meaning:} gradually become widely known or understood

\textbf{Example:} The findings took years to seep into public awareness.

\section{Synonyms and antonyms}
\sectionrule
\subsection{Close synonyms}
\subsubsection{leak}
\textbf{Part of speech:} verb

\textbf{Applies to:} Flow slowly through small openings

\textbf{Shared meaning:} Both can describe liquid or gas escaping unintentionally.

\textbf{Difference:} Leak emphasizes unintended escape through a hole or damaged container and may be fast or slow. Seep specifically emphasizes slow movement in small amounts, often through a surface.

\textbf{Example:} Oil leaked from the damaged tank.

\subsubsection{ooze}
\textbf{Part of speech:} verb

\textbf{Applies to:} Flow slowly through small openings

\textbf{Shared meaning:} Both describe very slow movement of liquid.

\textbf{Difference:} Ooze often describes a thick liquid moving slowly out of something. Seep commonly describes a thinner liquid passing gradually through a material or opening.

\textbf{Example:} Sap oozed from the damaged branch.

\subsubsection{permeate}
\textbf{Part of speech:} verb

\textbf{Applies to:} Spread or become known gradually

\textbf{Shared meaning:} Both can describe an influence gradually entering an area or situation.

\textbf{Difference:} Permeate means spread throughout something and focuses on the final wide presence. Seep focuses on the slow, gradual movement into it.

\textbf{Example:} A sense of optimism permeated the discussion.

\subsection{Direct or useful antonyms}
\subsubsection{gush}
\textbf{Part of speech:} verb

\textbf{Applies to:} Flow slowly through small openings

\textbf{Opposition:} Gush describes liquid flowing out suddenly and forcefully, whereas seep describes slow movement in small amounts.

\textbf{Example:} Water gushed from the broken main.

\section{Words learners may confuse}
\sectionrule
\subsection{leak}
\textbf{Part of speech:} verb

\textbf{Why learners confuse them:} Both often describe unwanted movement of liquid through an opening.

\textbf{Key difference:} A container, pipe, or roof can leak, and the escaping substance can leak from it. With seep, the slowly moving liquid or gas is normally the subject.

\textbf{seep:} Water seeped through the wall.

\textbf{leak:} The damaged roof leaked during the storm.

\subsection{sip}
\textbf{Part of speech:} verb

\textbf{Why learners confuse them:} The spellings and pronunciations are similar, but the vowel sounds and meanings differ.

\textbf{Key difference:} Seep means move slowly through small openings; sip means drink a small amount at a time.

\textbf{seep:} Moisture seeped into the wood.

\textbf{sip:} She sipped the hot tea slowly.

\section{Etymology}
\sectionrule
\textbf{Origin:} Seep developed as a variant of the older dialect word sipe, meaning to drain or ooze.

\textbf{Memory link:} Imagine water silently easing through tiny pores: slow movement is the key idea in seep.

\section{Practice exercises}
\sectionrule
\subsection{Word family choice}
Choose the best member of the word family for each blank.

\begin{enumerate}
  \item Engineers found evidence of water \_\_\_\_\_ beneath the barrier.
  \item Chemicals may \_\_\_\_\_ into the soil after heavy rain.
\end{enumerate}

\subsection{Correct or incorrect?}
Decide whether each sentence is correct. Correct the inaccurate ones.

\begin{enumerate}
  \item Water seeped through the porous rock.
  \item The damaged pipe seeped water onto the floor.
  \item News of the discovery gradually seeped out.
  \item The flood seeped rapidly through the city streets.
\end{enumerate}

\subsection{Collocation practice}
Complete each sentence with a natural collocation.

\begin{enumerate}
  \item Rainwater began to seep \_\_\_\_\_ the cracks in the roof.
  \item Confidential information slowly seeped \_\_\_\_\_ to the press.
  \item The chemical can seep \_\_\_\_\_ groundwater and remain there for years.
\end{enumerate}

\subsection{Complete answer key}
\subsubsection{Word family choice}
\begin{enumerate}
  \item \textbf{seepage.} Engineers found evidence of water seepage beneath the barrier. \emph{Use the noun seepage after evidence of.}
  \item \textbf{seep.} Chemicals may seep into the soil after heavy rain. \emph{Use the base verb seep after may.}
\end{enumerate}

\subsubsection{Correct or incorrect?}
\begin{enumerate}
  \item \textbf{Correct} \emph{Seep through correctly describes slow movement across a material.}
  \item \textbf{Incorrect}. Water seeped from the damaged pipe onto the floor. \emph{The liquid, not the pipe, is normally the subject of seep.}
  \item \textbf{Correct} \emph{Seep out can figuratively describe information becoming known little by little.}
  \item \textbf{Incorrect}. The flood swept rapidly through the city streets. \emph{Seep conflicts with rapid, forceful movement.}
\end{enumerate}

\subsubsection{Collocation practice}
\begin{enumerate}
  \item \textbf{through.} Rainwater began to seep through the cracks in the roof. \emph{Use through for the opening or material crossed.}
  \item \textbf{out.} Confidential information slowly seeped out to the press. \emph{Seep out describes information gradually becoming public.}
  \item \textbf{into.} The chemical can seep into groundwater and remain there for years. \emph{Use into for the destination.}
\end{enumerate}

\vfill
\begin{center}
\small\color{Muted} Created and compiled by \textbf{Amir Kooshky}\\
\href{https://t.me/KooshkyTOEFL}{Telegram Channel}\quad\textbullet\quad
\href{https://www.instagram.com/kooshkytoefl}{Instagram}\quad\textbullet\quad
\href{https://t.me/KooshkyTOEFL\_pv}{Personal Telegram}
\end{center}
\end{document}
